\chapter{Bộ lọc chính tả tiếng Việt cho khai thác neo}
\label{app:vi-filter}

Phụ lục này mô tả chi tiết bộ lọc chính tả được dùng ở Mục~\ref{subsec:anchor-mining}
để bảo đảm các ứng viên khai thác cặp dịch là từ tiếng Việt hợp lệ. Một token chỉ được
giữ lại khi vượt qua \textit{toàn bộ} các quy tắc sau:

\begin{enumerate}
  \item \textbf{Ký tự bản địa}: loại mọi token chứa \textit{f, j, w, z} -- các ký tự
        không thuộc bảng chữ cái tiếng Việt.
  \item \textbf{Độ dài}: tối đa 7 ký tự cho một âm tiết.
  \item \textbf{Một cụm nguyên âm}: token phải chứa \textit{đúng một} cụm nguyên âm
        liền mạch (kể cả nguyên âm mang dấu thanh).
  \item \textbf{Quy tắc phụ âm đầu}:
        \begin{itemize}
          \item \textit{k, gh, ngh} bắt buộc đứng trước các nguyên âm hàng trước
                \textit{i, e, ê, y};
          \item ngược lại, \textit{c, g, ng} không được đứng trước các nguyên âm này;
          \item \textit{q} luôn đi kèm \textit{u};
          \item \textit{p} không đứng một mình làm phụ âm đầu (loại các dạng vay mượn
                như \textit{pô, pin}).
        \end{itemize}
  \item \textbf{Tổ hợp nguyên âm ngoại lai}: loại các chuỗi không xuất hiện trong âm
        tiết tiếng Việt như \textit{ea, ee, oo, ii, uu, ion, eon}.
  \item \textbf{Nguyên âm chưa trọn vẹn}: các nguyên âm như \textit{ă, â} (và các tổ
        hợp \textit{iê, uô, ươ}) không được đứng cuối âm tiết mà phải có phụ âm cuối
        đi kèm (loại các mảnh từ như \textit{nhắ, tầ, nghiê}).
  \item \textbf{Phụ âm đầu/cuối hợp lệ}: phụ âm đầu phải thuộc tập
        \{$\varnothing$, b, c, ch, d, đ, g, gh, gi, h, k, kh, l, m, n, ng, ngh, nh,
        ph, q, qu, r, s, t, th, tr, v, x\}; phụ âm cuối phải thuộc tập
        \{$\varnothing$, c, ch, m, n, ng, nh, p, t\}.
  \item \textbf{Quy tắc thanh--phụ âm cuối}: âm tiết kết thúc bằng \textit{c, ch, p,
        t} bắt buộc mang thanh sắc hoặc thanh nặng (loại các dạng sai như
        \textit{sec, nhac, mit}).
  \item \textbf{Danh sách loại trừ}: một danh sách nhỏ các chuỗi thỏa cấu trúc âm
        tiết nhưng là nhiễu trong thực tế (ví dụ mảnh phiên âm, từ cổ ít dùng) được
        loại thủ công.
\end{enumerate}

Ngoài ra, ở Tầng 3 của quá trình khai thác cặp dịch (Mục~\ref{subsec:anchor-mining}),
một danh sách chặn (blacklist) gồm các từ chức năng và từ dễ nhập nhằng của tiếng Việt
(\textit{các, những, được, để, thế, cái,~\ldots}) được áp dụng trước khi dịch, nhằm
tránh tạo ra các cặp neo kém tin cậy từ những từ có phân bố ngữ nghĩa rộng.
